\section{Conclusion}
\label{sec:conclusion}

We introduced friction score (0-4: automated extraction + templates + validation + API) as a quantifiable metric capturing repository design quality, and demonstrated through cross-platform analysis of 10,000+ datasets that friction-reduction features correlate with metadata completeness. HuggingFace (friction=3) shows 55-66\% optional field presence versus UCI (friction=0) showing 0-15\%, with effect sizes spanning 45-61 percentage points (Cohen's $h=1.0$-1.8, $p<0.0001$).

Required fields (license, version) show 74-95\% presence across platforms with weaker friction effects (15-21pp HF-UCI difference, Cramér's $V=0.1$-0.2), validating enforcement as a distinct but coupled mechanism. Within-platform comparison confirms friction features reduce entry cost: API-uploaded datasets show 12.1pp higher completeness than manual uploads ($p<0.0001$, Cohen's $d=0.621$).

Our contributions include: (1) first cross-platform friction study at 10,000+ dataset scale, (2) mechanism validation distinguishing enforcement from friction reduction, (3) quantitative effect sizes with practical significance, and (4) repository design insights from gradient analysis.

\subsection{Future Directions}

\textbf{Randomized Controlled Trials:} Deploy feature ablation experiments on production platforms (e.g., A/B test templates vs no templates) to measure causal effects of individual friction features. Current gradient analysis provides decomposition estimates (templates est.\ 25-30pp, API est.\ 10-15pp) but requires experimental validation.

\textbf{Longitudinal Analysis:} Track metadata completeness over time to measure decay patterns, platform evolution effects, and persistence of friction-driven gains. Single-timepoint design cannot distinguish permanent gains from temporary compliance.

\textbf{Quality Assessment Beyond Presence:} Develop semantic analysis tools to measure metadata QUALITY (accuracy, completeness, usefulness) rather than just EXISTENCE. Binary presence detection misses low-quality placeholders passing parsing thresholds.

\textbf{Cross-Domain Validation:} Extend friction score framework to genomics (GEO, SRA), chemistry (PubChem), astronomy (SDSS) repositories to test generalizability beyond ML domain. Domain-specific metadata schemas may require friction feature adaptations.

\textbf{User Studies:} Measure friction effects across contributor experience levels (novice vs expert) and upload contexts (initial submission vs update) to identify equity considerations and targeted interventions.

\textbf{Automated Schema Alignment:} Develop ontology mapping tools to replace manual cross-platform schema normalization (currently 82.7\% semantic agreement). Automated alignment would enable larger-scale multi-platform studies.

\textbf{Weighted Friction Metrics:} Develop feature-weighted friction scores (e.g., templates=2, API=1, validation=0.5) based on validated marginal effects from ablation experiments. Current binary sum treats all features equally despite gradient evidence suggesting differential effects.

This work provides the first quantitative evidence linking friction reduction to metadata completeness at scale, establishing a foundation for data-driven repository design decisions and opening paths toward automated, friction-minimized scientific data sharing infrastructure.
